%==============================================================
% CHAPTER 1 - GENERAL INTRODUCTION
%==============================================================
\chapter{GENERAL INTRODUCTION}

\section{Introduction}

The healthcare landscape in Cameroon faces significant challenges characterised by resource limitations, uneven distribution of medical professionals, and systemic inefficiencies that compromise patient care quality and accessibility~\cite{minsante_2020}. Digital health solutions have emerged globally as transformative tools for bridging these gaps~\cite{who2019telemedicine,kvedar2014connected}. This project, \textbf{Clinix}, proposes a comprehensive mobile health platform designed specifically for the Cameroonian context. Clinix creates an integrated ecosystem connecting verified healthcare providers with patients through a secure, user-friendly digital interface, incorporating both traditional telemedicine features and AI-driven consultations.

\section{Background of the Study}

Cameroon's healthcare system is under sustained pressure. The country has roughly one physician per 5{,}000 inhabitants, well below the World Health Organization's recommended minimum, and that already-thin supply is unevenly distributed: most specialists are concentrated in a handful of cities while large rural and peri-urban populations are served by under-staffed district hospitals and informal home-care arrangements~\cite{minsante_2020}. Patients in regional centres routinely queue for hours, sometimes whole days, for consultations that could be resolved in fifteen minutes of focused triage. Laboratory turnaround is slow, prescription handling is largely paper-based, and medication adherence among patients with chronic conditions is widely reported as poor~\cite{wamala2013barriers}.

In parallel, Cameroon's digital infrastructure has matured rapidly. Mobile phone penetration exceeds eighty percent and smartphone adoption has grown sharply in urban areas~\cite{itu2023digital}. Mobile-money usage is now mainstream: MTN Mobile Money and Orange Money together process the bulk of small-value transactions, far exceeding card payments outside the formal banking sector~\cite{gsma2023mobilemoney}. This combination, widespread phones plus widespread mobile money, makes Cameroon a strong candidate for the kind of patient-facing digital-health platform that has already transformed primary care in the United States (Teladoc Health~\cite{teladoc}, K~Health~\cite{k_health}) and Rwanda (babyl), yet no platform has been built natively for the Cameroonian English-speaking market with these local realities at its core.

The professional supply side is similarly underserved. Many newly graduated doctors, nurses, midwives, and laboratory technicians remain underemployed within the formal hospital system and operate informally through social media advertisements and word-of-mouth home-care services~\cite{combi2016telemedicine}. There is presently no structured digital platform in Cameroon that verifies the credentials of these professionals, makes them discoverable to patients, processes payments transparently, and routes follow-up clinical artefacts (prescriptions, medical records, referrals) back to the patient. The result is a market that is simultaneously under-served on the patient side and under-utilised on the professional side, with no trusted layer in between.

Clinix is conceptualised against this background as a localised, secure, and scalable digital healthcare solution tailored to the Cameroonian English-speaking context, inspired by proven international telemedicine systems while addressing local healthcare realities head-on: low-bandwidth-tolerant design, mobile-money-native payments, mandatory provider verification cross-checked against the Cameroon Medical Council register, and an AI consultation layer that helps patients decide whether a hospital visit is even needed.

\section{Statement of the Problem}

The healthcare system in many parts of Cameroon is overwhelmed by increasing patient demand and limited healthcare infrastructure. Patients routinely experience long waiting times before consultations, delays in laboratory tests and prescription services, stress and frustration due to overcrowded hospitals, and difficulty locating nearby healthcare providers during emergencies.

At the same time, newly graduated healthcare professionals face limited employment opportunities within formal healthcare institutions. Many resort to offering home-based services without structured systems for authentication, scheduling, or patient trust, which leaves vulnerable patients exposed to unverified practitioners.

The absence of a centralised, secure, and intelligent digital platform that connects patients directly with verified healthcare providers has created a significant gap in healthcare delivery. There is also a lack of AI-assisted preliminary consultation tools that could help patients make informed decisions before visiting healthcare facilities. This project addresses these challenges by developing \textbf{Clinix}, a comprehensive mobile health platform that integrates telemedicine, AI consultation, provider verification, and GPS-based healthcare discovery.

\section{Objectives of the Study}

This section states the general objective of the project and the four specific objectives that operationalise it.

\subsection{General Objective}

To design, develop, and implement a comprehensive mobile health platform (Clinix) that digitally connects verified healthcare providers with patients in the English-speaking regions of Cameroon, integrating AI-assisted consultations, secure communication channels, and location-based services to improve healthcare accessibility, efficiency, and quality.

\subsection{Specific Objectives}

\begin{enumerate}
  \item To design a multi-tier architecture supporting mobile applications for both patients and healthcare providers, backed by a single REST API and a moderated administrative dashboard.
  \item To develop an AI-powered consultation module capable of symptom analysis, multimodal input, and ranked recommendation of verified providers tailored to the patient.
  \item To implement a robust provider-verification workflow in which submitted credentials are cross-checked against the Cameroon Medical Council register before a provider becomes visible to patients.
  \item To integrate location-based services, mobile-money payments, and bilingual-ready interfaces so the platform is usable on real Cameroonian networks, devices, and payment habits.
\end{enumerate}

\section{Significance of the Study}

The platform delivers concrete benefits to each stakeholder group. \textbf{Patients} gain reduced waiting times, increased access to verified healthcare professionals, the convenience of remote consultations, and a transparent way to discover nearby facilities. \textbf{Healthcare providers} gain expanded professional opportunities, flexible working arrangements, and digital tools for managing their patient base and earnings. The wider \textbf{healthcare system} benefits from reduced hospital congestion and improved service-delivery efficiency. \textbf{Policymakers} gain an evidence-based model for digital health integration and a potential framework for regulating telemedicine services. Finally, \textbf{researchers and developers} gain a documented contribution to academic knowledge in African mobile health systems and AI-driven healthcare solutions.

\section{Scope of the Study}

This study focuses on developing the Clinix platform for the English-speaking regions of Cameroon. The geographical focus is the Anglophone urban centres of \textbf{Buea}, \textbf{Bamenda}, \textbf{Limbe}, and \textbf{Kumba}, with plans for eventual expansion to other regions and full bilingual support in later releases. The target users are healthcare providers (doctors, nurses, midwives, laboratory technicians) and adult patients aged eighteen and above. Technically, the project covers an Android and iOS mobile application, a web-based administrative dashboard, and a cloud-hosted backend infrastructure. Functionally, the platform delivers virtual consultations, appointment scheduling, AI symptom checking, provider verification, mobile-money payment, medication reminders, home-care and lab-test bookings, and basic health-record management.

\section{Definition of Working Terms}

The following terms recur throughout this document and are defined here for reference. \textbf{Telemedicine} is the remote delivery of healthcare services using telecommunications technology. \textbf{mHealth} (mobile health) is medical and public-health practice supported by mobile devices. \textbf{Healthcare provider verification} is the process of authenticating professional credentials, licences, and qualifications. \textbf{AI consultation chatbot} refers to an artificial-intelligence system that simulates clinical conversation for preliminary medical guidance. \textbf{GPS-enabled map} is a digital map integrating Global Positioning System data for location services. \textbf{Superadmin} is the highest-level administrator with full system access and verification authority. \textbf{Mobile money} refers to electronic wallet services operated by telecom providers, principally MTN Mobile Money and Orange Money. \textbf{WebRTC} is a browser- and device-level standard for real-time audio, video, and data communication.

\section{Organisation of the Study}

This document is organised into five chapters. Chapter~I presents the introduction, problem statement, and research objectives. Chapter~II reviews relevant literature and existing telemedicine platforms. Chapter~III details the materials and methodology used for system development, including system analysis and system design. Chapter~IV covers implementation, results, and testing of the Clinix platform across its three sub-systems (mobile application, backend API, and administrative dashboard). Chapter~V presents conclusions, discusses technical limitations encountered during implementation, and outlines perspectives for further study.
